\section{Adjoint Equation}

To derive first-order optimality conditions, we introduce an adjoint variable associated with 
the stochastic state equation. The adjoint equation plays a central role in the computation of 
gradients and in the characterization of optimal controls.

\subsection{Lagrangian Functional}

Let \(y\) denote the state associated with a control \(u\). We introduce the Lagrangian functional

\begin{equation}
\begin{aligned}
\mathcal L(y,u,p)
=
&
\frac{1}{2}
\E
\left[
\int_0^T
|y-y_d|_H^2
,dt
\right]
+
\frac{\alpha}{2}
\E
\left[
\int_0^T
|u|_H^2
,dt
\right]
\
&
+
\E
\left[
\int_0^T
\left(
\frac{\partial y}{\partial t}
+
Ay
-
u,
, p
\right)_H
dt
\right],
\end{aligned}
\label{eq:lagrangian}
\end{equation}

where \(p\) is the adjoint variable.

For simplicity, the derivation is performed formally. A rigorous justification can be obtained 
using standard arguments from infinite-dimensional optimization theory.

\subsection{Variation with Respect to the State}

Let \(z\) be an admissible perturbation of the state variable. Differentiating the Lagrangian 
\eqref{eq:lagrangian} with respect to \(y\) in the direction \(z\) yields

\begin{equation}
\begin{aligned}
D_y\mathcal L(y,u,p)(z)
=
&
\E
\left[
\int_0^T
(y-y_d,z)_H
,dt
\right]
\
&
+
\E
\left[
\int_0^T
\left(
\frac{\partial z}{\partial t},
p
\right)_H
dt
\right]
\
&
+
\E
\left[
\int_0^T
(Az,p)_H
,dt
\right].
\end{aligned}
\label{eq:variation_y}
\end{equation}

We now integrate the time derivative term by parts:

\begin{equation}
\int_0^T
\left(
\frac{\partial z}{\partial t},
p
\right)_H
dt
=
(z(T),p(T))_H
-
(z(0),p(0))_H
-
\int_0^T
\left(
z,
\frac{\partial p}{\partial t}
\right)_H
dt.
\label{eq:integration_parts}
\end{equation}

Since the initial condition is fixed, we have

\[
z(0)=0.
\]

To eliminate the terminal contribution, we impose the final condition

\begin{equation}
p(T)=0.
\label{eq:adjoint_final}
\end{equation}

Substituting \eqref{eq:integration_parts} into \eqref{eq:variation_y} yields

\begin{equation}
\begin{aligned}
D_y\mathcal L(y,u,p)(z)
=
\E
\Bigg[
\int_0^T
\Big(
y-y_d
-
\frac{\partial p}{\partial t}
+
A^\ast p,
,z
\Big)_H
dt
\Bigg],
\end{aligned}
\label{eq:variation_final}
\end{equation}

where \(A^\ast\) denotes the adjoint operator of \(A\).

Since \(z\) is arbitrary, we obtain the adjoint equation.

\subsection{Adjoint Problem}

The adjoint variable \(p\) satisfies the backward evolution equation

\begin{equation}
-\frac{\partial p}{\partial t}
+
A^\ast p
=
y-y_d,
\qquad
\text{in }
D\times(0,T),
\label{eq:adjoint_abstract}
\end{equation}

together with the terminal condition

\begin{equation}
p(T)=0.
\label{eq:adjoint_terminal}
\end{equation}

For the advection--diffusion operator

\[
Ay
=
-\kappa \Delta y
+
\mathbf b\cdot\nabla y,
\]

the formal adjoint operator is given by

\begin{equation}
A^\ast p
=
-\kappa \Delta p
-
\nabla\cdot(\mathbf b\,p).
\label{eq:adjoint_operator}
\end{equation}

Consequently, the adjoint equation takes the form

\begin{equation}
-\frac{\partial p}{\partial t}
-\kappa \Delta p
-\nabla\cdot(\mathbf b\,p)
=
y-y_d,
\qquad
\text{in }
D\times(0,T).
\label{eq:adjoint_pde}
\end{equation}

The homogeneous Dirichlet boundary condition is

\begin{equation}
p=0,
\qquad
\text{on }
\partial D\times(0,T),
\label{eq:adjoint_bc}
\end{equation}

while the terminal condition is given by \eqref{eq:adjoint_terminal}.

The adjoint equation \eqref{eq:adjoint_pde} propagates information backward in time and provides 
the sensitivity of the cost functional with respect to perturbations of the state variable. 
It will be used in the next section to derive the first-order optimality conditions.
